% Build: cd docs/reports/reconstruction/methods && latexmk -xelatex pdc_reco_literature_and_air_ms_gap_20260426_zh.tex
\documentclass[a4paper, 11pt]{article}

\usepackage[margin=2.4cm]{geometry}
\usepackage{ctex}            % 中文支持
\usepackage{amsmath,amssymb}
\usepackage{booktabs}
\usepackage{multirow}
\usepackage{siunitx}
\usepackage{xcolor}
\usepackage{hyperref}
\usepackage{enumitem}
\usepackage{float}
\usepackage{graphicx}
\usepackage{tikz}
\usetikzlibrary{arrows.meta,positioning,calc}
\usepackage{longtable}
\usepackage{tabularx}

\definecolor{goodgreen}{RGB}{34,139,34}
\definecolor{warnred}{RGB}{200,50,50}

\hypersetup{
  colorlinks=true,
  linkcolor=blue!60!black,
  citecolor=blue!60!black,
  urlcolor=blue!60!black,
}

\title{\textbf{PDC 重建文献调研 \\ 与 SAMURAI 终端 air 多重散射 gap 分析}}
\author{TBT \\ RIKEN 仁科中心 / DPOL 实验}
\date{2026 年 4 月 26 日}

\begin{document}
\maketitle
\tableofcontents
\clearpage

\section{摘要 TL;DR}
% [content from Task 2]

\section{问题：SAMURAI 终端的 air 多重散射}
% [content from Task 3]
\subsection{质子飞行路径几何}
\subsection{实测 \texorpdfstring{$\sigma$}{sigma} 数值}
\subsection{与设计指标的差距}
\subsection{为什么这会限制科学产出}

\section{PDC 重建参考框架（文献）}
% [content from Tasks 4--6]
\subsection{硬件基础}
\subsection{算法主干}
\subsection{SAMURAI 各次实验中达到的性能}

\section{Gap 分析}
% [content from Tasks 7--8]
\subsection{性能：\texorpdfstring{$\sigma$}{sigma} 数值对比}
\subsection{方法学}
\subsection{物质预算与真空配置}
\subsection{实验配置与几何}

\section{扩展上下文：MINOS+SEASTAR 链}
% [content from Task 9]
\subsection{SEASTAR 计划时间线}
\subsection{配套探测器}
\subsection{综述与资源}
\subsection{广义文献中的注意事项}

\section{推论与下一步}
% [content from Task 10]
\subsection{短期数据获取行动项}
\subsection{中期模拟与分析方向}
\subsection{开放问题}
\subsection{Non-goals}

\section{参考文献}
% [content from Task 11]
\subsection{硬件与方法学（同行评议、基础）}
\subsection{算法参考}
\subsection{含 PDC 方法学的 PhD 论文}
\subsection{使用 PDC1/PDC2 的物理论文（确认/隐含）}
\subsection{排除/仅作上下文}
\subsection{综述与专题合集}
\subsection{建造提案与内部技术注释}
\subsection{项目内部参考}

\end{document}
